% !PRINTED 584
\weberpage{584}

Um die biquadratische Gleichung zu bilden, deren Wurzeln $\eta$,
$\eta_1$, $\eta_2$, $\eta_3$ sind, suchen wir zunächst die quadratische Gleichung
mit den Wurzeln $\eta$, $\eta_2$ und erhalten sie aus

\[
\begin{aligned}
2(\eta+\eta_2)&=-1+\sqrt{n}\\
2(\eta-\eta_2)&=(i,\eta)+(-i,\eta).
\end{aligned}
\]

Quadriren wir diese beiden Gleichungen und subtrahiren die
zweite von der ersten, so folgt wegen (30), (31) und (33)

\weq{35}{16\,\eta\,\eta_2=1+n-2\sqrt{n}-2n(-1)^f-2a(-1)^f\sqrt{n}.}

Durch (29) und (35) sind aber die Coëfficienten der gesuchten
quadratischen Gleichung bestimmt. Sie lautet:
\[
\eta^2+\frac{\eta}{2}+\frac{1+n-2n(-1)^f}{16}=\left(\frac{\eta}{2}+\frac{a(-1)^f+1}{8}\right)\sqrt{n},
\]
und die biquadratische Gleichung für die vier Perioden erhält
man, wenn man beiderseits quadrirt:

\weq{36}{\left(\eta^2+\frac{\eta}{2}+\frac{1+n-2n(-1)^f}{16}\right)^2-n\left(\frac{\eta}{2}+\frac{a(-1)^f+1}{8}\right)^2=0.}

Suchen wir den Coëfficienten der ersten Potenz von $\eta$, der
eine ganze Zahl sein muss, so finden wir dafür

\weq{37}{\frac{1+n-2n(-1)^f}{16}-n\,\frac{a(-1)^f+1}{8};}

$f$ ist gerade, wenn $n\equiv 1$, ungerade, wenn $n\equiv 5$ (mod 8) ist,
demnach ist
\[
1+n-2n(-1)^f\equiv 0\ (\mathrm{mod}\ 8).
\]
Daraus aber folgt wegen (37), dass $a(-1)^f+1$ durch 4
theilbar sein muss, also

\weq{38}{a\equiv -(-1)^f\ (\mathrm{mod}\ 4),}

und daraus ist nach (34) weiter zu schliessen, dass $b$ gerade
sein muss, also

\weq{39}{b\equiv 0\ (\mathrm{mod}\ 2).}

Die biquadratische Gleichung nimmt eine einfachere Gestalt
an, wenn wir

\weq{40}{4\eta+1=\xi}

setzen. Sie erhält dann nach (36) die Gestalt:

\weq{41}{[\xi^2+n(1-2(-1)^f)]^2-4n[\xi+(-1)^f a]^2=0,}

oder wenn wir die beiden Fälle $n\equiv 1$ oder $\equiv 5$ (mod 8) trennen:

\weq{42}{\begin{aligned}(\xi^2-n)^2-4n(\xi+a)^2&=0\\(\xi^2+3n)^2-4n(\xi-a)^2&=0.\end{aligned}}
